雅化的催化剂与风险都应落到模型变量，而不是停留在“锂价波动”或“行业竞争”的泛泛描述。最重要的事件是正式 H1：它既可能提升牛情景概率，也可能揭示利润与现金、库存或分部经济之间的缺口。

\section{催化剂日历：先验证，再重估}

\begin{exhibitbox}[表 9-1：未来验证事件与行动后果]
\begin{tabularx}{\exhibitboxwidth}{L{3.1cm}X X}
\toprule
事件 & 支持上修的可观测结果 & 行动后果 \\
\midrule
正式 2026H1 报告 & 归母落在预告区间，收入/毛利与锂盐改善相符，现金质量改善 & 复核基准情景，必要时提高牛情景概率 \\
H1 现金与营运资本披露 & OCF、应收、存货和库存变动有可复算解释 & 降低现金折价，允许重新评估估值倍数 \\
锂盐运营信息 & 销量、实现价格、成本或等效分部信息支持 H2 模型 & 允许将单位经济从假设转为证据 \\
民爆分部与项目披露 & 历史利润池未以新增应收/项目占款换取增长 & 提高防御性利润池的信用质量 \\
\bottomrule
\end{tabularx}
\end{exhibitbox}

\section{风险矩阵：现金和 H2 持续性优先于主题风险}

\begin{exhibitbox}[表 9-2：风险、概率与证伪阈值]
\begin{tabularx}{\exhibitboxwidth}{L{2.8cm}L{1.4cm}L{1.4cm}X X}
\toprule
风险 & 概率 & 影响 & 触发证据 & 模型/行动后果 \\
\midrule
利润未转现金 & 中高 & 高 & OCF 继续显著为负且应收/存货无解释 & 现金折价扩大，转 watchlist only \\
锂盐价差或销量回落 & 中高 & 高 & H2 收入/毛利不支持 9 万吨、10.5 万元/吨、32\% 毛利的基准组合 & 熊情景概率提高，EPS 与目标下调 \\
资源成本优势不兑现 & 中 & 中高 & 实际供给、自供比例或成本仍未披露/低于预期 & 保持零信用，禁止倍数上修 \\
民爆/矿服回款扰动 & 中 & 中 & 分部利润增长伴随项目应收或现金恶化 & 防御性利润池折现下调 \\
市场情绪反转 & 中 & 中 & 高波动延续且基本面验证不足 & 不以技术反弹加仓，维持减持纪律 \\
\bottomrule
\end{tabularx}
\end{exhibitbox}

\section{仓位与复核纪律}

当前行动是减持/避免新增仓位，而非基于短期价格回撤进行反向交易。若 H1 显示利润、现金和单位经济三者同时闭合，先将动作调整为中性并重估目标；若只验证利润而现金仍弱，维持防守；若预告兑现不足或营运资本继续显著吞噬现金，则退出为 `watchlist only`。牛情景并非不可实现，但必须由可观察的经营证据获得概率，而不是由客户、资源或行业叙事预支。
